\section{Testovanie výsledkov}

\subsection{Metodológia a testovacie dáta}

Štandardné hodnotenie odporúčacích systémov predpokladá prístup k reálnym používateľom a ich interakciám --
metriky ako Precision, recall či RSME sú zmysluplné iba vtedy
ak existujú overiteľné reálne preferencie, voči ktorým možno výstupy systému porovnávať \cite{ko2022survey}.
Keďže aplikácia ešte nebola nasadená pre reálnych používateľov a zber skutočnej spätnej väzby nebol v rámci tejto práce možný,
testovanie je realizované na synteticky generovaných dátach, ktoré simulujú správanie rôznych typov používateľov s rôznou históriou návštev.


Testovanie sme rozdelili na dve časti:
\begin{enumerate}[noitemsep, nolistsep]
    \item \textbf{Analýza vplyvu váh:} Sledovanie zmien odporúčaní pri systematickej
    zmene váhového vektora pre jedného používateľa a overenie citlivosti systému
    na jednotlivé faktory.
    \item \textbf{Overenie funkčnosti predpovede:} Overenie, či systém správne
    predpovedá relevantné udalosti na základe histórie používateľa, a nájdenie
    optimálnej konfigurácie váh pomocou metriky Precision@5.
\end{enumerate}


\subsection{Analýza vplyvu váh}
Cieľom tejto časti bolo pochopiť, ako jednotlivé faktory váhového vektora reálne
ovplyvňujú výsledné skóre odporúčaní, a identifikovať ktoré faktory sú pre
personalizáciu skutočne relevantné.

\subsubsection{Testovací používateľ}
Na účely testovania používame náhodného syntetického používateľa, ktorého budeme
nazývať Janko. Janko navštívil celkovo 8 udalostí: 4 športové, 3 vzdelávacie
a 1 z oblasti Business \& Tech. Konkrétne navštívené podujatia sú uvedené
v tabuľke~\ref{tab:konkretne_eventy}.

\begin{table}[H]
\centering
\caption{Navštívené udalosti a ich kategorizácia}
\label{tab:konkretne_eventy}
\small
\begin{tabular}{|l|l|}
\hline
\textbf{Podkategória}    & \textbf{Kategória} \\ \hline
Pitch Night              & Business \& Tech \\
Workshop (ručné práce)   & Education \\
Osobný rozvoj            & Education \\
Strength \& Conditioning & Sports \\
Team sports              & Sports \\
Winter sports            & Sports \\
Osobný rozvoj            & Education \\
Recreational sports      & Sports \\ \hline
\end{tabular}
\end{table}

Okrem histórie navštívených podujatí systém berie do úvahy aj udalosti označené
ako obľúbené. Tieto majú vyššiu váhu vďaka faktoru $B_{fav}$. Janko označil ako
obľúbenú udalosť v podkategórii \textit{Strength \& Conditioning} (Sports).

\subsubsection{Vizuálna analýza profilu}

Na obrázku\ref{fig:porovnanie_grafov} vidíme, že zatiaľ čo na úrovni hlavných
kategórií dominuje šport (50\,\%), na úrovni podkategórií sú záujmy veľmi
diverzifikované - každá podkategória má zastúpenie len 12,5–25\,\%.

\begin{figure}[H]
\centering
\begin{minipage}{0.48\textwidth}
    \centering
    \begin{tikzpicture}[scale=0.7]
        \tikzset{every node/.style={font=\scriptsize}}
        \pie[text=legend, radius=1.8,
             color={blue!60, green!60, orange!60}]{
            50/Sports,
            37.5/Education,
            12.5/Business \& Tech
        }
    \end{tikzpicture}
    \caption{Hlavné kategórie}
\end{minipage}
\hfill
\begin{minipage}{0.48\textwidth}
    \centering
    \begin{tikzpicture}[scale=0.7]
        \tikzset{every node/.style={font=\scriptsize}}
        \pie[text=legend, radius=1.8,
             color={cyan!60, yellow!60, red!60, gray!60, violet!60, brown!60, green!60}]{
            25/Osobný rozvoj,
            12.5/Pitch Night,
            12.5/Workshop,
            12.5/Strength \& Cond.,
            12.5/Team sports,
            12.5/Winter sports,
            12.5/Recreational sports
        }
    \end{tikzpicture}
    \caption{Podkategórie}
\end{minipage}
\caption{Preferencie používateľa Janka na úrovni kategórií a podkategórií}
\label{fig:porovnanie_grafov}
\end{figure}

\subsubsection{Výsledky testovania váh}

Pre výpočet finálneho skóre používame vážený súčet faktorov
(vzťah~\ref{eq:score_formula}):

\begin{equation}
Score = w_1 \cdot C_{main} + w_2 \cdot C_{sub} + w_3 \cdot D + w_4 \cdot R + w_5 \cdot P + w_6 \cdot B_{fav}
\label{eq:score_formula}
\end{equation}

Definovali sme sedem testovacích scenárov so systematickou zmenou 
váh $w_1$ až $w_6$,. Každý test obsahuje konfiguráciu váh,
  prehľad piatich najlepších (TOP 5) odporúčaní a podrobný rozbor príspevkov jednotlivých faktorov
 k výslednému skóre (detailné výstupy sú uvedené v Prílohe~\ref{priloha:vysledky}).

\paragraph{Výsledky pri rovnomernom rozdelení váh}

Pri rovnomernom rozdelení ($w_i \approx 16.67\,\%$) dominujú v odporúčaniach
udalosti s vysokým ratingom a obľúbené udalosti. Dôvodom sú vlastnosti
normalizovaných vstupov:

\begin{itemize}[noitemsep, nolistsep]
    \item \textbf{Obľúbený bonus:} Binárna hodnota (0 alebo 1) - automaticky
    posúva obľúbené udalosti o celú váhovú kategóriu vyššie.
    \item \textbf{Rating:} Väčšina podujatí má hodnotenie blízke 5 hviezdičkám,
    takže $R \approx 1.0$ - vytvára konštantný silný tlak na skóre.
    \item \textbf{Popularita:} Rastie pri napĺňaní kapacity; plne obsadené
    podujatia systém vyraďuje.
    \item \textbf{Vzdialenosť:} Stredne silný vplyv závislý od hustoty podujatí
    v okolí.
    \item \textbf{Hlavná kategória:} Kvôli diverzifikovanej histórii je $C_{main}$
    nízke. Pri rovnomernej váhe ho prekrývajú faktory s prirodzene vyššími hodnotami.
    \item \textbf{Podkategória:} Ešte nižší vplyv ako hlavná kategória -
    $C_{sub}$ je rozptýlené naprieč mnohými podkategóriami, takže pri rovnomernej
    váhe prakticky nevplýva na výsledné poradie.
\end{itemize}

Záver: čím viac chceme personalizovať podľa histórie, tým agresívnejšie musíme
zvyšovať váhy kategórií.

\paragraph{Porovnanie vplyvu faktorov pri dominantnej váhe 60\,\%}

\begin{table}[H]
\centering
\small
\begin{tabular}{|l|c|p{8cm}|}
\hline
\textbf{Faktor} & \textbf{Reálny vplyv*} & \textbf{Zdôvodnenie} \\
\hline
Obľúbený bonus   & 76\,\% & Binárny charakter (0/1) - najprudšie zmeny v poradí. \\
\hline
Rating           & 74\,\% & $R \approx 0.95$ - takmer fixný prírastok ku skóre. \\
\hline
Vzdialenosť      & 64\,\% & Vysoký rozptyl medzi blízkymi (0.85) a vzdialenými (0.35). \\
\hline
Hlavná kategória & 41\,\% & $C_{main} = 0.44$ pre šport (4/8 eventov). Závisí od histórie. \\
\hline
Podkategória     & 30\,\% & $C_{sub} = 0.11$ (1/8 eventov v každej podkategórii). Nízky
vplyv kvôli rozptýlenej histórii. \\
\hline
Popularita       & 31\,\% & Najslabší faktor - nízka obsadenosť v testovacej sade. \\
\hline
\end{tabular}
\caption{Vplyv faktorov pri 60\,\% konfigurácii (*priemerný podiel na skóre TOP~5)}
\label{tab:vplyv_faktorov}
\end{table}

\paragraph{Zhodnotenie faktora podkategórie}

Podkategória má najnižší reálny vplyv, jej zachovanie v algoritme má však
strategický význam:
\begin{itemize}[noitemsep, nolistsep]
    \item \textbf{Implicitné učenie:} Systém zachytí zmenu záujmov skôr, než si
    používateľ manuálne upraví profil.
    \item \textbf{Jemnosť odporúčaní:} Zabraňuje tomu, aby systém odporúčal všetky
    športové udalosti, keď používateľ preferuje konkrétne len jogu.
\end{itemize}

% ─────────────────────────────────────────────────────────────────────────────
\subsection{Overenie funkčnosti predpovede}
% ─────────────────────────────────────────────────────────────────────────────

\subsubsection{Cieľ}

Táto časť má dva ciele. Prvým je overiť, či systém vôbec dokáže správne predpovedať
relevantné udalosti na základe histórie - teda či základná funkčnosť funguje.
Druhým cieľom je nájsť optimálnu konfiguráciu váh pre hlavnú kategóriu a podkategóriu,
keďže z Časti~1 vyplynulo, že práve tieto dva faktory sú relevantné pre
personalizáciu.

\subsubsection{Overenie základnej funkčnosti}

Overovali sme, či systém správne predpovedá udalosti zodpovedajúce dominantnej
kategórii používateľa. Vytvorili sme 5~profilov s jasne definovanou dominantnou
kategóriou:

\begin{itemize}[noitemsep, nolistsep]
    \item \textbf{Sporťák} - 6$\times$ Sports, 1$\times$ Art \& Culture, 1$\times$ Food \& Drinks
    \item \textbf{Hudobník} - 5$\times$ Music \& Entertainment, 1$\times$ Sports
    \item \textbf{Umelec} - 4$\times$ Art \& Culture, 1$\times$ Music \& Entertainment
    \item \textbf{Technik} - 3$\times$ Business \& Tech, 1$\times$ Education
    \item \textbf{Foodie} - 4$\times$ Food \& Drinks, 1$\times$ Community \& Social
\end{itemize}

Predpoveď je úspešná (\textit{PASS}), ak aspoň 3~z TOP~5 odporúčaní patria do
dominantnej kategórie používateľa. Pre každý profil sme spustili algoritmus
s 10~konfiguráciami váh, kde sme menili pomer hlavná kategória / podkategória
od 100\,\%/0\,\% po 0\,\%/100\,\%:

\begin{table}[H]
\centering
\begin{tabular}{|l|c|c|}
\hline
\textbf{Konfigurácia} & \textbf{Hlavná kategória} & \textbf{Podkategória} \\
\hline
W01 & 100\,\% & 0\,\%   \\
W02 & 90\,\%  & 10\,\%  \\
W03 & 80\,\%  & 20\,\%  \\
W04 & 70\,\%  & 30\,\%  \\
W05 & 60\,\%  & 40\,\%  \\
W06 & 50\,\%  & 50\,\%  \\
W07 & 40\,\%  & 60\,\%  \\
W08 & 30\,\%  & 70\,\%  \\
W09 & 20\,\%  & 80\,\%  \\
W10 & 0\,\%   & 100\,\% \\
\hline
\end{tabular}
\caption{Konfigurácie váh pre overenie základnej funkčnosti}
\label{tab:w_configs}
\end{table}

Systém úspešne prešiel všetkými konfiguráciami pre všetkých 5~profilov - vo
všetkých prípadoch dokázal na základe histórie správne predpovedať relevantné
udalosti. Tento výsledok potvrdzuje, že content-based filtrovanie efektívne
extrahuje preferencie z histórie a aplikuje ich pri generovaní odporúčaní.

\subsubsection{Testovanie citlivosti predpovede}

Po potvrdení základnej funkčnosti sme hľadali optimálnu konfiguráciu pomeru
hlavná kategória / podkategória. Testovanie prebehlo na 50~syntetických
používateľoch s počtom navštívených udalostí od 1 do 50 a datasete 150~náhodných
udalostí. Pre každého používateľa sme získali referenčný TOP~5 na základe plnej
histórie, následne sme skrývali každý navštívený event a sledovali, či sa TOP~5
zmenil. Zmena TOP~5 po skrytí eventu je pozitívny signál - systém reaguje na
aktuálnu históriu a prispôsobuje odporúčania.

\begin{table}[H]
\centering
\begin{tabular}{|l|c|c|c|}
\hline
\textbf{Konfigurácia} & \textbf{Hlavná kat.} & \textbf{Podkategória} & \textbf{AVG citlivosť} \\
\hline
T01 & 100\,\% & 0\,\%   & 14\,\% \\
T02 & 90\,\%  & 10\,\%  & 35\,\% \\
T03 & 80\,\%  & 20\,\%  & 35\,\% \\
T04 & 70\,\%  & 30\,\%  & 36\,\% \\
T05 & 60\,\%  & 40\,\%  & 36\,\% \\
T06 & 50\,\%  & 50\,\%  & 40\,\% \\
T07 & 40\,\%  & 60\,\%  & 35\,\% \\
T08 & 30\,\%  & 70\,\%  & 35\,\% \\
T09 & 20\,\%  & 80\,\%  & 34\,\% \\
T10 & 0\,\%   & 100\,\% & 34\,\% \\
\hline
\end{tabular}
\caption{Výsledky testovania citlivosti - podiel prípadov, kde sa TOP~5 zmenil po skrytí eventu}
\label{tab:sensitivity}
\end{table}

T01 (100\,\% hlavná kategória) dosiahla najnižšiu citlivosť 14\,\% - systém
reagoval na zmenu histórie len v 14\,\% prípadov. Konfigurácia T06 (50\,\%/50\,\%)
dosiahla maximum 40\,\%, pričom T10 (100\,\% podkategória) skončila na 34\,\%.
Kombinácia oboch faktorov v rovnakom pomere teda reaguje na zmeny histórie lepšie
ako ktorýkoľvek faktor samostatne.

\subsubsection{Obmedzenia metodiky}

Leave-one-out metodika má na syntetických dátach niekoľko prirodzených obmedzení.
Používatelia s menej ako 5~navštívenými udalosťami predstavujú cold start scenár,
kde história po skrytí jednej udalosti nestačí na spoľahlivú personalizáciu.
Pri diverzifikovanej histórii môže skrytie udalosti zmeniť dominantnú kategóriu
profilu - systém sa správa správne, ale metrika to hodnotí ako miss. Ak používateľ
navštívil kategóriu len raz, po jej skrytí klesne $C_{main}$ na nulu - ide o
prirodzenú vlastnosť content-based filtra, nie o chybu implementácie.

\subsection{Záver testovania}

Testovanie potvrdilo, že systém je schopný na základe histórie správne predpovedať
relevantné udalosti - vo všetkých 5~profiloch a naprieč všetkými konfiguráciami
váh systém úspešne odporučil udalosti zodpovedajúce dominantným záujmom
používateľa. Na základe výsledkov testovania citlivosti sme zvolili konfiguráciu
T06 (50\,\%/50\,\%) ako optimálnu pre produkčné nasadenie. Treba však poznamenať,
že táto konfigurácia nemusí byť vhodná v každom prípade - pri používateľovi
s históriou obsahujúcou iba jeden typ udalosti v rámci kategórie môže systém
odporúčať správne kategórie, ale nemusí správne predpovedať konkrétny typ podujatia.
Kvalita predpovede závisí nielen od nastavenia váh, ale aj od hĺbky a diverzity
histórie používateľa.


